\chapter{\MakeUppercase{Methodology and Testing}}

\section{Methodology}
The development of Badal followed an agile methodology with iterative cycles of feature development and testing.

\subsection{Scanning Methodology}
To ensure thorough coverage while respecting API rate limits, the system employs:
\begin{itemize}
    \item \textbf{Paginators:} Using Boto3 paginators to handle large numbers of resources.
    \item \textbf{Targeted Discovery:} Focusing on high-impact services (Compute, Identity, Storage).
    \item \textbf{Package Extraction:} For EC2, the system uses AWS SSM (Systems Manager) to run non-intrusive commands and collect installed package versions. For Lambda, it analyzes layer configurations.
\end{itemize}

\subsection{AI Analysis Methodology}
The integration with Google Gemini uses a "Structured Prompting" technique. The prompt includes:
\begin{itemize}
    \item \textbf{Role Definition:} Instructing the LLM to act as a senior cloud security architect.
    \item \textbf{Context Provision:} Providing the full resource and dependency JSON report.
    \item \textbf{Constraint Enforcement:} Requiring the output to be a valid JSON with specific keys (priority, problem, solution, cli_command).
\end{itemize}

\section{Testing Strategy}
The system was tested using a multi-tiered approach:

\subsection{Unit Testing}
Individual components like the Risk Predictor and the NVD Parser were tested with mock data to ensure mathematical correctness and robust error handling.

\subsection{Integration Testing}
The integration between the FastAPI backend and the AWS scanners was tested using a dedicated AWS sandbox environment containing intentionally misconfigured resources (e.g., public S3 buckets, roles with AdministratorAccess).

\subsection{Validation of AI Solutions}
AI-generated CLI commands were manually verified for syntax and effectiveness. The system's ability to prioritize critical vulnerabilities was benchmarked against manual security audits of the same environment.

\section{Performance Analysis}
\begin{table}[h]
\centering
\caption{System Performance Metrics}
\renewcommand{\arraystretch}{1.3}
\begin{tabular}{|l|l|l|}
\hline
\textbf{Task} & \textbf{Environment Size} & \textbf{Average Time} \\
\hline
Resource Discovery & 50 Resources & 12 Seconds \\
Dependency Mapping & 50 Resources & 2 Seconds \\
Vulnerability Scan (NVD) & 20 Packages & 45 Seconds (due to rate limits) \\
AI Analysis (Gemini) & Full Report & 8 Seconds \\
\hline
\end{tabular}
\end{table}

\section{Security Considerations}
During testing, specific attention was paid to the secure handling of credentials.
\begin{itemize}
    \item Credentials are never stored on the server disk.
    \item JWT tokens are used to maintain sessions, with an expiry of 1 hour.
    \item All external API calls (AWS, NVD, Gemini) are performed over encrypted HTTPS channels.
\end{itemize}
